\documentclass[a4paper, 10pt, conference]{ieeeconf}
\usepackage[utf8]{inputenc}
\usepackage{graphicx}
\usepackage{amsmath}
\usepackage{amssymb}
\usepackage{url}
\usepackage{xcolor}
\usepackage{hyperref}
%
% ##########################
% include as many packages you need here
% ##########################
%

%
% ##########################
% Do not change this header
% ##########################
%
\title{WACV 2027 UPAR Challenge\\ \textit{Fact sheet}} 
\begin{document}
\maketitle
This is the fact sheet's template for the \href{https://vap.aau.dk/rws/challenge/}{WACV 2027 UPAR Challenge}. Please fill out the following sections carefully in a scientific writing style. Then, send the compressed project (in \texttt{.zip} format), i.e., the generated PDF, \texttt{.tex}, \texttt{.bib} and any additional files to \url{upar.challenge@gmail.com}.
%
% end of the header
%------------------
%

% ##########################
% Next, you are asked to fill in the required information 
% (please, do not remove/change any item)
% ##########################

\section{Team details}\label{teamdetails}
\begin{itemize}
    \item Team leader name: Kaveh Safavigerdini
    \item Username(s) on Codabench: kavehsfv-957256-957256\_VID-Docker-GPU
    \item Team leader affiliation: Electrical Engineering and Computer Science Department, University of Missouri - Columbia, Mo, USA
    % \item Team leader address: 
    % \item Team leader phone number: 
    \item Team leader email: ksgh2@umsystem.edu
    \item Name of other team members (and affiliation): ---
    \item Team website URL (if any): ---
    \item Competition track (mark with X one single option)\footnote{If you participated in more than one competition track, you need to share with the organizers one fact sheet per track.}:
    \begin{itemize}
        \item (X) Track 1: Pedestrian Attribute Recognition.
        \item ( ) Track 2: Attribute-based Person Retrieval.
        \item ( ) Track 3: Human Pose Estimation.
    \end{itemize}
\end{itemize}
%
% end of 'Team details'
%------------------
%

% ##########################
% Next, you are asked to fill in the required information 
% Here, you are free to add figures and tables, and expand the default sections to better explain your method and contributions.
% ##########################

\section{Contribution details}
%

\subsection{Title of the contribution}
Fine-Tuning a CLIP ViT-B/16 Visual Encoder with Tempered Class-Balanced BCE and Coordinate-Ascent Logit Calibration for Unified Pedestrian Attribute Recognition.

\subsection{Representative image / workflow diagram of the method}
A pedestrian crop is letterboxed to $224\times224$ with grey ($128$) padding and CLIP-normalized, then encoded by the fine-tuned CLIP ViT-B/16 visual tower~\cite{radford2021clip} ($12$ blocks, width $768$, QuickGELU, class-token pooling, $512$-d projection) and a linear head ($\mathbb{R}^{512}\to\mathbb{R}^{40}$) into raw attribute logits $\mathbf{z}\in\mathbb{R}^{40}$. Adding a validation-calibrated bias $\mathbf{b}\in\mathbb{R}^{40}$ yields final probabilities $\hat{\mathbf{p}}=\sigma(\mathbf{z}+\mathbf{b})$ for the challenge's fixed $0.5$ threshold.

\subsection{Detailed method description}
\textbf{Architecture.} We adopt the visual tower of OpenAI CLIP ViT-B/16 ($16\times16$ patches, $12$ pre-norm Transformer blocks of width $768$ with $12$ heads and QuickGELU, class-token pooling, and a $768\times512$ projection) initialized from public \texttt{openai} weights~\cite{radford2021clip}, topped with a single linear layer ($\mathbb{R}^{512}\to\mathbb{R}^{40}$) for the $40$ UPAR attributes~\cite{Specker_2023_WACV,cormier2024upar}. The submission is implemented in pure PyTorch and Pillow with no external dependencies.

\textbf{Training objective.} To counter extreme class imbalance on the training split (positive rates $p_a$ span $0.17\%$ for \texttt{LowerBody-Color-Purple} to $93.47\%$ for \texttt{Age-Adult}), we use binary cross-entropy with tempered positive-class weights:
\begin{equation}
w_a = \left(\frac{1-p_a}{\max(p_a,\,10^{-6})}\right)^{0.5}, \qquad a \in \{1,\dots,40\},
\end{equation}
compressing weights to $[0.26,\,24.16]$ to boost rare-attribute recall without gradient instability.

\textbf{Optimization \& data.} We train exclusively on the UPAR training split ($97{,}669$ images: $10{,}000$ Market-1501, $79{,}001$ PA-100K, $8{,}668$ PETA) for $6$ epochs with batch size $256$ (\texttt{drop\_last=True}), AdamW (weight decay $0.05$), differential learning rates ($2\times10^{-5}$ backbone, $1\times10^{-3}$ head), $500$-step linear warmup plus cosine decay, $\text{bf16}$ AMP, and gradient-norm clipping at $1.0$. The validation split ($33{,}407$ images) is used solely to select the best checkpoint (epoch $2$, raw validation $\text{Challenge\_Avg}=0.8751$) and calibrate logit biases.

\textbf{Post-processing (coordinate-ascent logit calibration).} Since evaluation enforces a fixed $0.5$ threshold while ranking by $\text{Challenge\_Avg}$ (harmonic mean of $\text{mA}$ and $\text{Instance\_F1}$), we tune a $40$-d logit bias $\mathbf{b}\in[-3,3]^{40}$ (and temperature $T\in\{1,2\}$; $T=1$ optimal) on the validation set via $2$ coordinate-ascent sweeps over a $17$-point grid per attribute (step $0.375$). Predicting $\sigma(\mathbf{z}+\mathbf{b})$ implements per-attribute thresholds $\tau_a=\sigma(-b_a)$ under the $0.5$ rule ($32/40$ biases non-zero), lifting validation $\text{Challenge\_Avg}$ from $0.8751$ to $0.8835$ ($\text{mA}$: $0.8717\to0.8941$; $\text{Instance\_F1}$: $0.8784\to0.8731$).

\subsection{Challenge results}
Validation ($33{,}407$ images) and official test leaderboard ($36{,}343$ images) results:

\begin{table}[h]
\centering
\begin{tabular}{l c c c}
\hline
Metric & Val (raw) & Val (cal.) & \textbf{Test (leaderboard)} \\
\hline
\textbf{Challenge\_Avg} & 0.8751 & 0.8835 & \textbf{0.7599} \\
mA & 0.8717 & 0.8941 & 0.7716 \\
Label\_F1 & 0.7049 & 0.6673 & 0.5341 \\
Instance\_Acc & 0.8001 & 0.7913 & 0.6257 \\
Instance\_Prec & 0.8690 & 0.8494 & 0.7152 \\
Instance\_Rec & 0.8953 & 0.9070 & 0.8009 \\
Instance\_F1 & 0.8784 & 0.8731 & 0.7486 \\
\hline
\end{tabular}
\end{table}

\subsection{Final remarks}
\textbf{Pros.} (i) Fine-tuning a CLIP ViT-B/16 visual tower with a single linear head gives strong cross-domain features without part-localization branches. (ii) Tempered BCE weighting ($\gamma=0.5$) and coordinate-ascent logit calibration directly align predictions with the harmonic $\text{Challenge\_Avg}$ metric. (iii) Self-contained pure-PyTorch inference with batched $\text{bf16}$ execution.

\textbf{Cons.} (i) The validation-to-test drop ($0.8835\to0.7599$) reflects domain shift and mild bias-calibration overfitting to validation priors (trading $\text{Label\_F1}$ precision for $\text{mA}$ recall). (ii) No ensembling or test-time augmentation. (iii) Rare color attributes remain the main error source.

\section{Additional method details}
Please, reply if your challenge entry considered (or not) the following strategies and provide a brief explanation. For each question, mark with X one single option. \\

%
% mark with an X each question.
%

\begin{itemize}
%------------------------------------
\item \textbf{For the competition track associated with this fact sheet, you confirm that you have trained your model on the predefined datasets} (X)~Yes, (~)~No\\
Trained strictly on the predefined UPAR training split ($97{,}669$ images: Market-1501, PA-100K, PETA).
%------------------------------------

%------------------------------------
\item \textbf{Did you use any pre-trained model:} (X)~Yes, (~)~No.\\ If yes, please detail (e.g., model architecture, training strategy, train data, etc.).\\
%
Initialized from the OpenAI CLIP ViT-B/16 visual tower (\texttt{openai} weights via \texttt{open\_clip\_torch}, WIT-400M~\cite{radford2021clip}), fine-tuned end-to-end with a linear head ($\mathbb{R}^{512}\to\mathbb{R}^{40}$).
%------------------------------------

%------------------------------------
\item \textbf{Did you use external data?} (~)~Yes, (X)~No\\
If yes, please detail:\\
%
No external images, annotations, or pseudo-labels beyond the public CLIP initialization weights.
%------------------------------------

%------------------------------------
\item \textbf{Did you perform any data augmentation?} \\(X)~Yes, (~)~No\\
If yes, please detail:\\
%
Random horizontal flip ($p=0.5$) and tensor-space random erasing ($p=0.25$, rectangle size $\in[1/8,1/3)$ per side filled with $0.0$, the normalized CLIP mean).
%------------------------------------

%------------------------------------
\item \textbf{Did you use the provided validation set as part of your training set?} (~)~Yes, (X)~No\\
If yes, please detail:\\
%
Never used for gradient updates; used only for epoch-$2$ checkpoint selection and post-hoc logit-bias calibration.
%------------------------------------

%------------------------------------
\item \textbf{Did you use any regularization strategies/terms?} (X)~Yes, (~)~No\\
If yes, please detail:\\
%
AdamW weight decay ($0.05$), gradient $\ell_2$-norm clipping at $1.0$, random erasing ($p=0.25$), and early stopping at epoch $2$.
%------------------------------------

%------------------------------------
\item \textbf{Did you use any pre-processing steps?} (X)~Yes, (~)~No\\
If yes, please detail:\\
%
Aspect-ratio-preserving bilinear letterbox resize to $224\times224$ with grey ($128$) padding and CLIP normalization ($\mu=(0.4815, 0.4578, 0.4082)$, $\sigma=(0.2686, 0.2613, 0.2758)$).
%------------------------------------

%------------------------------------
\item \textbf{Did you use any post-processing steps?} (X)~Yes, (~)~No\\
If yes, please detail:\\
%
Per-attribute logit bias $\mathbf{b}\in[-3,3]^{40}$ tuned via $2$-sweep coordinate ascent ($17$-point grid, $T=1.0$) on validation $\text{Challenge\_Avg}$; inference outputs $\sigma(\mathbf{z}+\mathbf{b})$ thresholded at $0.5$.
%------------------------------------

%------------------------------------
\item \textbf{Did you use Vision Language Models (VLMs) of any kind?} (X)~Yes, (~)~No\\
If yes, please detail:\\
%
Only the visual tower of CLIP ViT-B/16~\cite{radford2021clip} is used as an image backbone; the text encoder and language supervision are unused.
%------------------------------------
\end{itemize}


\section{Code repository}
Link to a code repository with complete and detailed instructions so that the results obtained on Codabench can be reproduced locally. This includes a list of requirements, pre-trained models, and so on. Note, training code with instructions is also required. This is required for all participants to appear in the final ranking. \textbf{Organizers strongly encourage the use of docker to facilitate reproducibility}. \\

\textbf{Code repository:} \url{https://github.com/kavehsfv/UPAR_Track1}


\bibliographystyle{IEEEtran}
\bibliography{references}
\end{document}
